% =======================================================
% 3. CASI D'USO
% =======================================================
\section{Casi d'Uso}
[-]

\subsection{Attori}
\subsubsection{Attore primario 1: {[Nome]}}
{[Descrizione]}

\subsubsection{Attore primario 2: {[Nome]}}
{[Descrizione]}

\subsection{Diagrammi dei casi d'uso}

\subsubsection{UC1 - {[Nome Caso d'Uso]}}
% Inserisci qui l'immagine del diagramma UML
% \begin{figure}[ht]
	%    \centering
	%    \includegraphics[width=0.7\textwidth]{uml.png}
	%    \caption{Diagramma UML}
	% \end{figure}

\subsubsection{UC2 - {[Nome Caso d'Uso]}}
% \begin{figure}[ht] ... \end{figure}

\newpage
\subsection{Descrizione dettagliata casi d'uso}

% STRUTTURA TABELLARE
\subsubsection{UC1 - {[Titolo Caso d'Uso]}}
\begin{table}[ht]
	\centering
	\renewcommand{\arraystretch}{1.4}
	\begin{tabularx}{\textwidth}{|l|X|}
		\hline
		\textbf{Codice} & \textbf{UC1} \\
		\hline
		\textbf{Titolo} & {[-]} \\
		\hline
		\textbf{Attori} & {[-]} \\
		\hline
		\textbf{Descrizione} & {[-]} \\
		\hline
		\textbf{Precondizioni} & {[-]} \\
		\hline
		\textbf{Postcondizioni} & {[-]} \\
		\hline
		\textbf{Scenario Principale} & 
		\vspace{-0.4cm}
		\begin{enumerate}[leftmargin=*]
			\item [-]
			\item [-]
			\item [-]
		\end{enumerate} \\
		\hline
		\textbf{Scenari Alternativi} & 
		\vspace{-0.4cm}
		\begin{itemize}[leftmargin=*]
			\item [-]
		\end{itemize} \\
		\hline
		\textbf{Estensioni/Inclusioni} & 
		\vspace{-0.4cm}
		\begin{itemize}[leftmargin=*]
			\item \textbf [-]
			\item \textbf [-]
		\end{itemize} \\
		\hline
	\end{tabularx}
	\caption{UC1 - {[Titolo Caso d'Uso]}}
\end{table}

\vspace{0.5cm}
\subsubsection{UC1.1 - {[Titolo Sotto-Caso d'Uso]}}
\begin{table}[ht]
	\centering
	\renewcommand{\arraystretch}{1.4}
	\begin{tabularx}{\textwidth}{|l|X|}
		\hline
		\textbf{Codice} & \textbf{UC1.1} \\
		\hline
		\textbf{Titolo} & {[-]} \\
		\hline
		\textbf{Attori} & {[-]} \\
		\hline
		\textbf{Descrizione} & {[-]} \\
		\hline
		\textbf{Precondizioni} & {[-]} \\
		\hline
		\textbf{Postcondizioni} & {[-]} \\
		\hline
		\textbf{Scenario Principale} & 
		\vspace{-0.4cm}
		\begin{enumerate}[leftmargin=*]
			\item [-]
		\end{enumerate} \\
		\hline
	\end{tabularx}
	\caption{UC1.1 - {[Titolo Sotto-Caso d'Uso]}}
\end{table}